\clearpage\fancyhead[L]{\small COMET.}\hypertarget{native-347-lower}{}\label{native-347-lower}
\begin{center}\textbf{COMET.}\par\vspace{4pt}\hrule\vspace{4pt}\fontsize{10}{12}\selectfont Göttingische gelehrte Anzeigen. Stück 201. Seite 2001\,..\,2004. 1811 Decemb. 19.\par\vspace{4pt}\hrule\vspace{4pt}\end{center}

\par\indent Im 151 Stück dieser Blätten [S. 337 d. B.] haben wir die ersten hiesigen Beobachtungen des grossen diesjährigen Cometen, nebst den ersten verbesserten parabolischen Elementen, angezeigt. Letztere waren nur als vorläufige angegeben; sie fingen bald an, sich einige Minuten von den Beobachtungen zu entfernen, und zu Anfang October war dieser Unterschied bereits auf einen Viertelsgrad angewachsen. Eine leichte Verbesserung der parabolischen Elemente reichte indessen hin, diesen Fehler wegzuschaffen: diese zweiten verbesserten Elemente sind im Octoberheft der Monatl. Correspondenz [S. 338 d. B.] bekannt gemacht. Anfangs November war der Fehler dieser Elemente, bei deren Berechnung nur einige einzelne Beobachtungen zum Grunde gelegen hatten, wieder auf einige Minuten angewachsen: und Hr. Prof. \textsc{Gauss} hielt es daher für interessant, zu untersuchen, in wie fern man dies schon als einen Beweis von Ellipticität der Bahn ansehen könne. Er liess die hierzu nöthigen Rechnungen unter seiner Aufsicht von Hrn. \textsc{Nicolai} ausführen, von dessen ausgezeichneter Geschicklichkeit und Sorgfalt im astronomischen Calcul wir schon früher in diesen Blättern Proben mitgetheilt haben. Es wurden der grösste Theil der sämmtlichen Beobachtungen des Hrn. von \textsc{Zach}, in der ersten Periode der Sichtbarkeit des Cometen, und eine grosse Anzahl neuerer Beobachtungen, die bis zum 6. November reichten, zum Grunde gelegt, und aus deren Vergleichung mit den letzten Elementen des Hrn. Prof. \textsc{Gauss} vier Normalörter abgeleitet. Das Resultat war, dass mit Hülfe einer nur sehr kleinen Correction der letztern Elemente, die neuern Beobachtungen sich genau darstellen liessen, während bei den ältern Beobachtungen nur kleine Differenzen zurück blieben, nemlich $16''$ in der Länge, und $28''$ in der Breite bei dem ersten Normalorte vom 16. April, und $51''$ in Länge, und $120''$ in der Breite den 18. Mai bei dem zweiten. Obgleich nicht anzunehmen ist, dass der letztere Normalort, das Mittel aus einer grossen Anzahl freilich nicht sehr genauer Beobachtungen, wirklich ganz mit einem so grossen Fehler behaftet sei, so ist derselbe doch noch zu klein, um bei der Ungewissheit, ein wie grosser Theil davon noch auf Rechnung des Normalorts selbst zu setzen sei, eine einigermassen zuverlässige Bestimmung der Ellipse gründen zu können. Daher hielt Hr. Prof. Gauss es für besser, dieses Geschäft noch zu verschieben, bis spätere Beobachtungen Etwas zu entscheiden in den Stand setzen. Auf alle Fälle ist die Umlaufszeit weit über 1000 Jahre. Die verbesserten parabolischen Elemente nach Hrn. \textsc{Nicolai}'s Rechnung sind folgende:

\Needspace{18\baselineskip}
\begin{center}\fontsize{10}{12}\selectfont\setlength{\tabcolsep}{4pt}\renewcommand{\arraystretch}{1.15}\begin{tabular}{l}
Durchgang durch die Sonnennähe 1811. 12. Sept.\\
\quad $6^{\mathrm u}\;30^{\mathrm m}\;55^{\mathrm s}$ Mittl. Z. in Göttingen\\
\end{tabular}\end{center}
\begin{center}\fontsize{10}{12}\selectfont\setlength{\tabcolsep}{3pt}\renewcommand{\arraystretch}{1.15}\begin{tabular}{llr}
Logarithm. des kleinsten Abstandes &  & $0.0151048$\\
Länge der Sonnennähe & .\;.\;.\;.\;. & $75^\circ\;1'\;44''3$\\
Länge des aufsteig. Knoten & .\;.\;. & $140\;21\;57.5$\\
\end{tabular}\end{center}
\noindent Beide siderisch ruhend, und von der Nachtgleiche des 12. September gezählt.
\begin{center}\fontsize{10}{12}\selectfont\setlength{\tabcolsep}{3pt}\renewcommand{\arraystretch}{1.15}\begin{tabular}{llr}
Neigung der Bahn & .\;.\;.\;.\;.\;.\;. & $73^\circ\;4'\;30''9$\\
\end{tabular}\end{center}
\noindent Bewegung rückläufig.

\SourceLeaf{348}
\par\indent Zur Erleichterung der Beobachtungen im Januar des nächsten Jahres hat Hr. \textsc{Nicolai} noch folgende Ephemeride berechnet, welche wir ganz hiehersetzen, wenn gleich die Sichtbarkeit des Cometen früher aufhören wird. \mbox{--- --- ---}


\par\indent Vorstehendes war bereits zum Abdruck niedergeschrieben, als Hr. Prof. \textsc{Gauss} in einem Schreiben des Hrn. v. \textsc{Lindenau} die Nachricht erhielt, dass Hr. \textsc{Pons} in Marseille den 16. November noch einen neuen Cometen im Eridanus entdeckt habe. Folgende Beobachtungen des Hrn. von \textsc{Zach} zu Marseille waren dem Schreiben beigefügt:

\begin{center}
\fontsize{10}{12}\selectfont\setlength{\tabcolsep}{3pt}\renewcommand{\arraystretch}{1.15}\begin{tabular}{l|c|c|c}
1811 & \shortstack{Mittlere Zeit\\in Marseille} & \shortstack{Gerade\\Aufsteigung} & \shortstack{Südliche\\Abweichung} \\ \hline
Novemb. 17 & $10^{\mathrm u}\;23^{\mathrm m}$ & $67^\circ.25'$ & $25^\circ\,52'$ \\
18 & $11\;\;11\;\;17^{\mathrm s}3$ & $67\;\;14\;\;39''8$ & $25\;\;24\;\;\;8''6$ \\
19 & $\;9\;\;59\;\;32.8$ & $67\;\;\;4\;\;59.6$ & $24\;\;54\;\;\;8.5$ \\
20 & $10\;\;\;8\;\;37.7$ & $.66\;\;56\;\;\;8.2$ & $24\;\;18\;\;\;9.2$ \\
21 & $10\;\;14\;\;45.5$ & $66\;\;46\;\;53.0$ & $23\;\;41\;\;47.8$ \\
\end{tabular}
\end{center}


\Needspace{22\baselineskip}\par\indent Gleich am Abend des 9. Decembers, wo diese Nachricht eingegangen war, begünstigte ein sehr heiterer Himmel die Aufsuchung des Cometen auf hiesiger Sternwarte. Hr. Prof. \textsc{Harding} nahm ihn auch sofort mit einem Cometensucher in der Nähe eines Sterns siebenter Grösse wahr, mit welchem Hr. Prof. \textsc{Gauss} ihn mehrere Male am Kreismicrometer verglich. Die scheinbare Position des Sterns wurde aus der \emph{Histoire céleste} zu $64^\circ\,11'\,23''1$ gerader Aufsteigung, $10^\circ\,22'\,2''5$ südl. Abweichung bestimmt, woraus folgende Position des Cometen sich ergab:

\Needspace{9\baselineskip}
\begin{quote}\raggedright 1811 Dec. 9. $10^{\mathrm u}\;6^{\mathrm m}\;52^{\mathrm s}$ ger. Aufst. $63^\circ\;49'\;41''4$, südliche Abweichung $10^\circ\;21'\;55''5$\par\medskip Am 11. Dec. $10^{\mathrm u}\;34^{\mathrm m}\;1^{\mathrm s}$ M. Z. wurde gefunden: ger. Aufst. $63^\circ\;33'\;20''5$, südl. Abw. $8^\circ\;39'\;54''7$\end{quote}


\par\indent Dieser neue Comet ist übrigens bis jetzt sehr klein und lichtschwach, etwa so hell wie ein Stern achter Grösse; er scheint indess, da seine Bewegung sich beschleunigt, der Erde näher zu kommen; auch lässt er sich, weil er einen bestimmtern Mittelpunkt zeigt, als der grosse, am Kreismicrometer besser beobachten.


\par\indent Den 12. December.


